\paragraph{The specification and the SDKs}
SEP-2145 is close to settled before anyone votes on it. SEP-1303 already
reclassified input validation failures as tool execution errors, arguing that a model
can only correct an error it sees~\cite{sep1303}, and SEP-2145 proposes the same for
unknown tools~\cite{sep2145}. Adopting it would not change how servers behave. It
would change the document to describe what \ErrAsResultRate{} of responding servers
already do. Rejecting it would make \ErrAsResultCount{} servers wrong
overnight, and nobody is going to fix them. Nor is there a coherent minority to
protect: of the servers that do not answer with \texttt{isError},
\UnknownProtoErrCount{} emit the protocol error the specification illustrates,
\UnknownWrongCodeCount{} emit a different code, and \UnknownProseCount{} report the
failure only in prose.

The second point matters more than the first. This behavior was not chosen. Server
authors did not weigh the two mechanisms and pick one; they inherited a default that
is the same across the official TypeScript and Python SDKs and the third-party
wrappers built on them. So \ErrAsResultRate{} is not a vote, and it should not be
read as community consensus about the right design. It is a count of who used an SDK.
The same fact cuts the other way, and more usefully: one change at the library layer
moves thousands of servers at once, in whichever direction the specification chooses.
The limit of that lever is visible here too. The unknown-tool fix recorded on the
TypeScript SDK's tracker in March 2026 shipped four months before our census, and the
rate did not move.

\paragraph{Silent execution is a different kind of finding}
The unknown-tool divergence is a disagreement about which channel carries an error.
Either way the client is told that something failed. Silent execution tells the
client nothing failed. The specification obliges every server to validate the inputs
it is given, no pending proposal would relax that obligation, and so there is no
argument that the specification will catch up.

A client that receives one of these answers cannot tell it from a correct one. There
is no error code, no \texttt{isError}, nothing to branch on. A tool that scans code
for package imports, handed the integer \texttt{12345}, replied \texttt{CLEAN}, which
is what a real clean scan looks like. The agent cannot retry, cannot repair the
argument, and cannot warn the user. It carries the wrong result forward. And
\NSilentExec{} is a lower bound, because we credited any error-like wording in a reply
as a report. This is also not \NSilentExec{} careless authors. \SilentExecTS{} of them
are built on the official TypeScript SDK, whose low-level \texttt{Server} API
publishes a tool's \texttt{inputSchema} and then never checks an incoming call against
it (Section~\ref{sec:consequences}). The default does the publishing and leaves out
the checking.

\paragraph{Startup failure and what a registry could do about it}
Two in five installable servers never serve the protocol, and nothing in a registry
listing says so. A registry could try a handshake once before listing an entry and
show the outcome. It need not delist anything; it need only tell people. The check is
cheap, and we ran it across the whole ecosystem.

Two further gaps in the registry record are ours to report because we hit them. The
registry does not record an entry point, which left \NEntrypointBlocked{} PyPI servers
impossible to launch from registry data alone, and cost our own measurement about two
percentage points of apparent PyPI yield until we corrected for it
(\PypiRateRaw{} uncorrected against \PypiRateCorr{} after). And required configuration is not recorded either: \NNeedsConfig{}
servers failed for no reason other than a setting that nothing told us to supply.
Neither is a defect in the server. Both are things a client cannot learn without
running the software.

\paragraph{What running the servers shows that reading them cannot}
None of the above is visible to a study that reads code or metadata. The static
ecosystem studies~\cite{mcpfirstglance,mcplandscape} cannot see a server that does not
start, and they cannot see one that answers a poisoned call as though it were
ordinary. The artifact's presence in a registry says almost nothing about the
artifact. The same lesson is reported for Jupyter notebooks, where publishing one and
running it turned out to be different
questions~\cite{notebookreproducibility}, and for Docker Hub images that share a
name~\cite{dockerhubimages}. It is also a caution about registries as a research
source~\cite{huggingfacehub}: our own entry-point problem was exactly that, a
launcher reading registry fields and getting them wrong.

One contrast in our data sharpens the point. Runnability differs sharply between the
two package registries, \NpmRate{} against \PypiRateCorr{}, while conformance among
the servers that do run does not. Packaging problems are per-ecosystem and
per-package. Protocol behavior is per-SDK. Fixing conformance means changing a few
libraries. Fixing runnability means changing thousands of packages. Those are very
different amounts of work, and only the first has a lever short enough to pull.

\paragraph{A moving specification}
The lifecycle we measure is being retired. Revision 2026-07-28 drops the
initialization handshake and negotiates a version on every request instead, with an
optional \texttt{server/discover} call for clients that want to ask up
front~\cite{mcpversioning}. Handshake-based behavior is now described as a backward
compatibility concern. Our population sits well behind that boundary:
\PctVerOffered{} of responders settled on 2025-06-18, and \NVerDowngraded{} answered
with something older still. This census records an installed base at the point its
mechanism left the specification, and measures the distance between the two. A server
published today inherits whichever default its SDK shipped with, and the revision that
removes the handshake reaches those servers only when their authors upgrade, which the
version spread here suggests is slow.

\paragraph{Recommendations}
For the specification: accept SEP-2145, state it with an RFC-2119 keyword rather than
an example, and add the unknown-tool case to the conformance suite so that it is
tested rather than described.

For SDK authors: validate \texttt{tools/call} arguments against the declared
\texttt{inputSchema} before dispatch, in the low-level \texttt{Server} API as well as
the high-level one, and check the tool name against the declared list. Most authors
never touch a default, so the default is the behavior.

For agent framework authors: do not treat the absence of \texttt{isError} as evidence
that a call succeeded. Validate arguments against the declared schema on the client
side, because \SilentExecRate{} of responding servers do not, and treat a server as
unproven until it completes a handshake.

For registry maintainers: run a handshake at publish time and show the result, and
record the entry point and the required configuration in the listing.
